\documentclass[11pt]{article}
\usepackage[utf8]{inputenc}
\usepackage[T1]{fontenc}
\usepackage[a4paper,margin=2.5cm]{geometry}
\usepackage{amsmath,amssymb}
\usepackage{booktabs}
\usepackage{graphicx}
\usepackage{algorithm}
\usepackage{algorithmic}
\usepackage[round,authoryear]{natbib}
\usepackage{url}
\usepackage[hidelinks]{hyperref}
\usepackage{xcolor}

\newcommand{\todo}[1]{\textcolor{red}{[TODO: #1]}}

\title{Forget the distant past: training windows and window combination\\
for automatic ARIMA and exponential smoothing on long time series}
\author{Draft (2026-10-27) \todo{authors}}
\date{}

\begin{document}
\maketitle

\begin{abstract}
Statistical forecasting methods are usually fitted to the whole of the
available history. For long series, with thousands of observations, this
assumes one model with fixed parameters holds throughout, which is rarely
plausible. We study whether automatic ARIMA and exponential smoothing (ETS),
as implemented in the \texttt{aeon} toolkit, forecast better when fitted to a
window of recent data, and how to choose or combine windows automatically. On
LSSH-99, a collection of 99 long series of up to 10{,}000 observations, with
one-step rolling-origin evaluation, fixed windows of 300--1{,}000 observations
reduce the MAE of AutoARIMA and AutoETS by 6--8\% relative to the whole series.
Combining the forecasts from a geometric grid of windows is more reliable than
choosing one: the median over windows reduces error by 11--13\%, and weighting
windows by their recent out-of-sample error by 12--15\%, improving about 70\%
of series. On the 5{,}070 M4 series with at least 500
observations, under the standard M4 protocol, the gains are smaller (OWA
0.920 to 0.894--0.905 for \texttt{AutoARIMA}, 0.945 to 0.922--0.926 for
\texttt{AutoETS}), concentrated on hourly, monthly and weekly series, with no
effect on daily series. The median over windows, and the median of the half of
the windows with the lowest out-of-sample error, are close to the best scheme
on both benchmarks; weights from in-sample fit, Akaike-style weights and
selection of the best-fitting window are no better, or worse, on M4,
consistent with the forecast combination puzzle. We propose a simple wrapper
that turns any forecaster into a window-combining forecaster.
\end{abstract}

% =====================================================================
\section{Introduction}

Automatic statistical forecasters such as \texttt{auto.arima} and
\texttt{ets}~\citep{hyndman08auto,hyndman02state,hyndman08book} remain strong
benchmarks: in the M4 competition combinations of them were among the most
accurate methods~\citep{makridakis20m4}. They are almost always fitted to the
whole of the available history. For the short series that dominate the M
competitions this is sensible, since there is little data to spare. For long
series it is less obvious. A series of 10{,}000 hourly observations covers more
than a year; its level, seasonal profile and dynamics are unlikely to be
constant over that period, and a model estimated on all of it averages over
regimes that may no longer apply. Practical advice recognises this:
\citet{hyndman21fpp3} note that ``most time series models do not work well for
very long time series'' and suggest, for short-horizon forecasts, to ``throw
away the earliest observations and only fit a model to the most recent
observations''. They give no rule for how many to keep.

The question of how much history to use has a long history in econometrics,
where it is framed as the choice of an estimation window under structural
breaks~\citep{pesaran07window} and has led to methods that combine forecasts
across estimation windows~\citep{pesaran11averaging}. That work concerns
mostly regression and autoregressive models with a fixed specification applied
to macroeconomic and financial series. Here we ask the same question of
automatic ETS and ARIMA, which select a model on every fit, on benchmark
collections of long univariate series, and with the interaction between
windows and the seasonal adjustment used in forecasting practice.

We make three contributions.
\begin{enumerate}
\item An evaluation of training windows for automatic ARIMA and ETS on LSSH-99,
a collection of 99 long series~\citep{banwell26forecasting}, and on the long
series of M4, showing that fitting all the data is generally no better, and
often clearly worse, than fitting recent data.
\item A comparison of ways to choose a window or combine forecasts across
windows: per-series and rolling selection, means, medians, trimmed means,
weights from in-sample fit normalised for window length, Akaike-style weights,
and weights from recent out-of-sample error.
\item A simple forecaster wrapper derived from these results (its
implementation and evaluation as an ordinary forecaster are in progress), and a
store of all window forecasts that allows further combination schemes to be
evaluated without refitting.
\end{enumerate}

% =====================================================================
\section{Related work}
\label{sec:related}

\paragraph{Choosing the estimation window.}
\citet{pesaran07window} frame the choice as a trade-off between the bias from
using data from before a break and the variance from a shorter sample, and
propose cross-validation methods to choose a single window together with
methods that weight or pool forecasts from windows of different lengths.
\citet{inoue17rolling} choose the rolling window to minimise an estimate of the
forecaster's quadratic loss, by analogy with bandwidth selection in
nonparametric estimation. \citet{clark09combining} show that combining
forecasts from an expanding (recursive) and a rolling window is at least as
accurate as either. \citet{rossi12robust} make the related point for
evaluation: tests of predictive ability depend on the window, and repeating
them over windows invites data snooping.

\paragraph{Combining across windows.}
\citet{pesaran11averaging} average one model's forecasts over estimation
windows and find lower bias and RMSFE than any single window for all but the
smallest breaks. \citet{tian14breaks} weight windows by the estimated
probability that each point is the most recent break.
\citet{assenmacher08swiss} find averaging over windows at least as effective as
averaging over models for Swiss macroeconomic forecasts, with little gain from
weighting schemes over equal weights, and \citet{jungmittag16air} find
window-averaged SARIMA forecasts of monthly air travel generally better than
single windows. Most recently, \citet{petersen26windows} compare training
window selection methods for regression and random forests on non-stationary
psychological time series: a simple average across windows and a super learner
ensemble perform best, while selecting a single window by cross-validation
performs poorly, and the expanding window is best only when the coefficients
are stable.

\paragraph{Down-weighting instead of truncating.}
\citet{pesaran13optimal} derive the weights on past observations that minimise
mean squared forecast error under breaks, close to exponential under
continuous change; \citet{giraitis13adaptive} choose the discount rate of
rolling and exponentially weighted estimators by cross-validation. Outside
econometrics, adaptive windowing for drifting data streams
\citep{bifet07adwin} and dynamic ensembles that weight forecasters over time
\citep{cerqueira17arbitrated} address related problems.

\paragraph{Forecast combination.}
Equal weights are hard to beat: estimated weights rarely improve on the simple
average~\citep{stock04combination}, because of the error in estimating them
\citep{smith09puzzle,claeskens16puzzle}. Trimmed and Winsorised means are
slightly more accurate and more robust than the mean~\citep{jose08robust}.
\citet{kolassa11akaike} combines ETS models fitted to the same series with
Akaike weights, which beat selection by an information criterion. Bagging ETS
over bootstrapped series~\citep{bergmeir16bagging} improves on ETS, mainly
through the uncertainty in model selection~\citep{petropoulos18bagging}, a
mechanism that also operates when an automatic forecaster selects different
models on different windows.

\paragraph{What is new here.}
We have not found an evaluation of window selection or combination for
automatic ETS and ARIMA on benchmark collections of long series, of the median
across windows, or of the interaction between windows and seasonal adjustment;
nor a forecasting library that offers combination across estimation windows
as a component.

% =====================================================================
\section{Methods}
\label{sec:methods}

\subsection{Forecasters}
We use four forecasters from \texttt{aeon}~\citep{aeon24jmlr}, at the code
released in version 1.7: \texttt{AutoARIMA} (stepwise order search over
non-seasonal ARIMA models, fitted by conditional sum of squares),
\texttt{AutoETS} (selection over error, trend and seasonal components by AIC),
and two fixed models, \texttt{ARIMA}(1,0,1) and simple exponential smoothing
\texttt{ETS}(A,N,N), the toolkit defaults. As in the M4 benchmarks, the
non-seasonal \texttt{AutoARIMA} and \texttt{ARIMA} are fitted to series
deseasonalised by classical multiplicative decomposition when a 90\%
autocorrelation test finds seasonality, and the forecasts are reseasonalised;
\texttt{AutoETS} is given the seasonal period; \texttt{ETS}(A,N,N) is fitted to
the raw series.

\subsection{Training windows}
Given a series $y_1,\dots,y_t$ observed up to the forecast origin $t$, the
forecaster fitted to window $w$ uses $y_{t-w+1},\dots,y_t$; $w=0$ denotes the
whole series. We use the grid
\[
W=\{50, 100, 200, 300, 500, 750, 1000, 1500, 2000, 3000, 5000, 7500\}
\]
plus the whole series, skipping windows at least as long as the history. Let
$\hat y^{(w)}_{t+h}$ be the forecast from window $w$.

\subsection{Choosing a window}
\emph{Validation} selection chooses, per series, the window with the lowest
error on a block of origins before the test period (LSSH-99), or on the last
$h$ training observations (M4). \emph{Rolling} selection chooses, at each
origin, the window with the lowest mean one-step error over the previous 30
origins. \emph{Global} selection chooses one window for all series.

\subsection{Combining windows}
A combination forecasts $\sum_w \lambda_w \hat y^{(w)}_{t+h}$ with weights
$\lambda_w\ge 0$ summing to one, or a robust location of the window forecasts.
\begin{itemize}
\item \emph{Mean}, \emph{median}, and the \emph{trimmed mean} (20\% from each
end).
\item \emph{In-sample fit.} With $e^{(w)}_s$ the in-sample one-step residuals of
the fit to window $w$, we use per-point scores so that windows of different
lengths are comparable: the mean squared residual
$\mathrm{MSE}_w$, and the relative MAE $r_w$, the mean absolute residual divided
by the mean absolute one-step change over the same points (the in-sample error
of the naive forecast), which also removes the local scale of the series.
Weights are $\lambda_w\propto \mathrm{MSE}_w^{-1}$, $\propto r_w^{-1}$ or
$\propto r_w^{-2}$.
\item \emph{Akaike-style.} The AIC is not comparable across samples of
different size, so we scale the per-point Gaussian log-likelihood to a common
length $n_0=100$: $\lambda_w\propto\exp\{-\tfrac{n_0}{2}(\log\mathrm{MSE}_w-
\min_v\log\mathrm{MSE}_v)\}$.
\item \emph{Best in-sample window}: the window with the smallest $r_w$.
\item \emph{Out-of-sample error.} $\lambda_w\propto \mathrm{OOS}_w^{-1}$, with
$\mathrm{OOS}_w$ the mean squared one-step error of window $w$ at the previous 30
origins (LSSH-99) or on the M4 holdout; and the median of the better half of
windows by $\mathrm{OOS}_w$.
\item \emph{Minimum window.} The mean, median, inverse in-sample MSE and inverse
out-of-sample MSE restricted to windows of at least 300 observations (and the
whole series).
\end{itemize}
For multiplicative-error ETS models the residuals returned are relative errors;
we convert them to the scale of the data by multiplying by the fitted values.

% =====================================================================
\section{Experimental design}
\label{sec:design}

\paragraph{LSSH-99.}
The collection of~\citet{banwell26forecasting}: 99 long series (1{,}008 to
10{,}000 observations), most from the Monash archive~\citep{godahewa21monash},
including 57 long M4 series, with seasonal periods from the Monash conventions
(48 half-hourly, 24 hourly, 7 daily, 144 ten-minute, 1440 minutely, 52 weekly,
12 monthly). Each forecaster is refitted at each of the last 60 observations
and forecasts one step ahead; origins 1--30 form a validation period and
origins 31--60 the test period, the same test origins as in the evaluation of
\texttt{aeon}'s forecasters by~\citet{banwell26forecasting}. Per series, the
MAE over the test origins is divided by that of the naive (last value)
forecast; methods are compared by the geometric mean of this ratio over 98
series (one series is constant, and on one the naive forecast is exact), on
the 36 sub-daily series and the 62 daily-or-slower series. For each method
against the whole series we count series better and worse and apply a
Wilcoxon signed-rank test to the log ratios, Holm-corrected within each
forecaster. The whole-series forecasts are identical to those of the reference
runs.

\paragraph{M4 long series.}
The 5{,}070 M4 series~\citep{makridakis20m4} with at least 500 training
observations (2{,}946 daily, 1{,}403 monthly, 414 hourly, 252 weekly, 43
quarterly, 12 yearly; median length 1{,}430), with the standard protocol: one
origin at the end of the training series, the competition horizons, sMAPE and
MASE averaged over the forecast points, and OWA relative to the published
Naive2 forecasts on the same series. Each window is also fitted to the training
series without its last $h$ observations, forecasting them, to provide a
holdout for selection and for out-of-sample weights.

\paragraph{Reproducibility.}
All forecasts are stored per window, with the in-sample fit scores, so that any
further combination can be evaluated without refitting. Code, results and the
store are in the accompanying repository~\todo{URL}.

% =====================================================================
\section{Results}
\label{sec:results}

\subsection{Does a window beat the whole series?}

Table~\ref{tab:lssh-windows} gives the LSSH-99 relative MAE of each fixed
window. For \texttt{AutoARIMA}, \texttt{AutoETS} and \texttt{ARIMA}(1,0,1),
windows of 300--1{,}000 observations beat the whole series by 6--8\%; very
short windows (50--200) are no better than the whole series, and the error
rises back towards the whole-series value as the window grows beyond 2{,}000.
Simple exponential smoothing is almost unaffected (1--2\%). No single fixed
window is significantly better than the whole series after correction for the
13 windows tested.

\begin{table}[t]
\centering\small
\caption{LSSH-99: relative MAE (geometric mean over 98 series of MAE divided by
the naive MAE; below 1 is better than naive) of each forecaster fitted to the
last $w$ observations at each of 30 one-step origins. Bold: best per row.}
\label{tab:lssh-windows}
\resizebox{\textwidth}{!}{% generated by make_tables.py from results/aeon-8585c7bf; do not edit
\begin{tabular}{lccccccccccccc}
\toprule
 & \multicolumn{13}{c}{Training window (observations)} \\
\cmidrule(lr){2-14}
Forecaster & Whole & 50 & 100 & 200 & 300 & 500 & 750 & 1000 & 1500 & 2000 & 3000 & 5000 & 7500 \\
\midrule
\texttt{AutoARIMA} & 0.918 & 0.907 & 0.924 & 0.884 & 0.861 & \textbf{0.846} & 0.847 & 0.865 & 0.873 & 0.865 & 0.882 & 0.889 & 0.909 \\
\texttt{AutoETS} & 0.927 & 1.024 & 0.987 & 0.917 & 0.909 & 0.880 & 0.868 & \textbf{0.859} & 0.879 & 0.876 & 0.905 & 0.910 & 0.915 \\
\texttt{ARIMA}(1,0,1) & 0.932 & 0.929 & 0.946 & 0.936 & \textbf{0.862} & 0.866 & 0.864 & 0.866 & 0.875 & 0.882 & 0.914 & 0.918 & 0.931 \\
\texttt{ETS}(A,N,N) & 1.023 & 1.016 & 1.007 & 1.005 & \textbf{1.005} & 1.007 & 1.006 & 1.008 & 1.026 & 1.026 & 1.025 & 1.024 & 1.021 \\
\bottomrule
\end{tabular}
}
\end{table}

\subsection{Choosing or combining windows}

Table~\ref{tab:lssh-main} compares choosing a window with combining windows.
Choosing per series on a validation block beats the best fixed window for
\texttt{AutoARIMA} and \texttt{ARIMA}(1,0,1) but not for \texttt{AutoETS}; rolling selection, which re-chooses the window
at each origin from the previous 30 errors, gains 12--13\% for
\texttt{AutoARIMA} and \texttt{AutoETS}. Combination is at least as good and
more reliable: the median over windows improves on the whole series by 11.5\%
(\texttt{AutoARIMA}) and 13.2\% (\texttt{AutoETS}), for 66 and 73 of the 98
series, both significant, and it is better than the mean for every forecaster.
Weighting the windows by their recent out-of-sample error is better still for
\texttt{AutoETS} (0.789 against 0.927 for the whole series, 72 series better).
The oracle, which picks the best window per series with knowledge of the test
errors, shows that more is available (0.73--0.75).

\begin{table}[t]
\centering\small
\caption{LSSH-99: relative MAE of the whole series, the best fixed window (in
brackets), window selection and window combination. Oracle: best window per
series chosen on the test errors (not a method).}
\label{tab:lssh-main}
\resizebox{\textwidth}{!}{% generated by make_tables.py from results/aeon-8585c7bf; do not edit
\begin{tabular}{lccccccc|c}
\toprule
 & Whole & Best fixed & \multicolumn{2}{c}{Selection} & \multicolumn{3}{c|}{Combination} & \\
\cmidrule(lr){4-5}\cmidrule(lr){6-8}
Forecaster & series & window & validation & rolling & mean & median & inv.\ OOS MSE & Oracle \\
\midrule
\texttt{AutoARIMA} & 0.918 & 0.846 (500) & 0.813 & \textbf{0.800} & 0.834 & 0.813 & 0.803 & 0.745 \\
\texttt{AutoETS} & 0.927 & 0.859 (1000) & 0.861 & 0.808 & 0.818 & 0.804 & \textbf{0.789} & 0.734 \\
\texttt{ARIMA}(1,0,1) & 0.932 & 0.862 (300) & 0.824 & \textbf{0.823} & 0.881 & 0.846 & 0.840 & 0.787 \\
\texttt{ETS}(A,N,N) & 1.023 & \textbf{1.005} (300) & 1.013 & 1.006 & 1.014 & 1.010 & 1.012 & 0.990 \\
\bottomrule
\end{tabular}
}
\end{table}

\subsection{Weighting schemes}

Table~\ref{tab:lssh-schemes} gives every combination scheme. Three findings
stand out. First, the gains over the median are small: the best weighted
schemes improve on it by 1--2\%, as the forecast combination
puzzle~\citep{smith09puzzle,claeskens16puzzle} leads one to expect. Second,
the schemes that rely most on in-sample fit quality, selecting the window with
the best in-sample fit and Akaike-style weights (which concentrate almost all
weight on the best-fitting window), are the weakest weighted schemes for
\texttt{AutoARIMA} and \texttt{AutoETS}, and not significantly better than the
whole series for \texttt{AutoARIMA}: in sample, a model fitted to fewer points
fits them more closely. Third, inverse
in-sample MSE weights do well (0.799 for \texttt{AutoARIMA}, 0.790 for
\texttt{AutoETS}). They are far from selection: the largest weight is a median
of about 0.2 (equal weights give 0.08), and they tilt the combination mildly
towards shorter windows, over which the series has been calmer (weighted
geometric mean window 480 against 597 under equal weights for
\texttt{AutoARIMA}, 538 for \texttt{AutoETS}). Unlike the relative MAE, this
score is not scale-free, so it rewards windows covering low-volatility periods,
which on these data acts as a mild recency weighting.

\begin{table}[t]
\centering\small
\caption{LSSH-99: relative MAE of each combination scheme, with the number of
series better / worse than the whole series. $^*$: significantly different from
the whole series (Wilcoxon signed-rank, Holm-corrected, $p<0.05$). Bold: best
per forecaster.}
\label{tab:lssh-schemes}
\resizebox{\textwidth}{!}{% generated by make_tables.py from results/aeon-8585c7bf; do not edit
\begin{tabular}{lcccc}
\toprule
Combination & \texttt{AutoARIMA} & \texttt{AutoETS} & \texttt{ARIMA}(1,0,1) & \texttt{ETS}(A,N,N) \\
\midrule
Whole series & 0.918 & 0.927 & 0.932 & 1.023 \\
Mean & 0.834$^*$ (66/32) & 0.818$^*$ (68/30) & 0.881 (57/41) & 1.014 (38/47) \\
Median & 0.813$^*$ (66/32) & 0.804$^*$ (73/25) & 0.846$^*$ (60/38) & 1.010 (27/24) \\
Trimmed mean (20\%) & 0.824$^*$ (68/30) & 0.806$^*$ (74/24) & 0.870$^*$ (58/40) & 1.012 (38/35) \\
Inverse in-sample MSE & \textbf{0.799}$^*$ (62/36) & 0.790$^*$ (72/26) & 0.844 (59/39) & 1.010 (35/55) \\
Inverse in-sample relative MAE & 0.815$^*$ (69/29) & 0.803$^*$ (69/29) & 0.850$^*$ (59/39) & 1.013 (39/46) \\
Inverse squared in-sample relative MAE & 0.808$^*$ (70/28) & 0.799$^*$ (71/27) & 0.845$^*$ (59/39) & 1.012 (38/46) \\
Akaike-style weights ($n_0=100$) & 0.825 (49/49) & 0.831$^*$ (61/37) & 0.852 (58/40) & 1.007 (37/54) \\
Best in-sample window & 0.836 (47/51) & 0.820$^*$ (59/39) & 0.851 (55/43) & 1.010 (27/33) \\
Inverse out-of-sample MSE & 0.803$^*$ (69/29) & \textbf{0.789}$^*$ (72/26) & 0.840$^*$ (59/39) & 1.012 (39/53) \\
Median of better half & 0.805$^*$ (65/33) & 0.796$^*$ (73/25) & \textbf{0.837}$^*$ (59/39) & \textbf{1.005} (27/28) \\
Mean, $w\geq300$ & 0.844$^*$ (63/35) & 0.842$^*$ (67/31) & 0.880$^*$ (61/37) & 1.016 (43/29) \\
Median, $w\geq300$ & 0.836$^*$ (63/35) & 0.837$^*$ (74/24) & 0.866$^*$ (60/38) & 1.019 (30/20) \\
Inverse in-sample MSE, $w\geq300$ & 0.835$^*$ (64/34) & 0.830$^*$ (71/27) & 0.863$^*$ (62/36) & 1.013 (44/46) \\
Inverse out-of-sample MSE, $w\geq300$ & 0.832$^*$ (64/34) & 0.827$^*$ (69/29) & 0.862$^*$ (61/37) & 1.014 (48/41) \\
\bottomrule
\end{tabular}
}
\end{table}

\subsection{Where windows help}

Table~\ref{tab:lssh-periods} breaks the LSSH-99 results down by sampling
frequency. The gains are concentrated on sub-daily data: hourly series
(\texttt{AutoARIMA} 0.774 to 0.543 with the median) and minutely wind power,
where the whole-series fits are worse than naive (1.49 and 1.19) and windows
bring them to 0.64--0.75. \texttt{AutoETS} also gains strongly on half-hourly
series (0.741 to 0.512), \texttt{AutoARIMA} does not. On weekly and monthly
series windows help moderately; on daily series, which are close to random
walks, they make little difference. A likely mechanism for the hourly series is that deseasonalisation
and the seasonal models estimate the daily profile from the window, so the
profile is local rather than averaged over a year.

\begin{table}[t]
\centering\small
\caption{LSSH-99 by sampling frequency (seasonal period): relative MAE of the
whole series, the median over windows and inverse out-of-sample MSE weights.}
\label{tab:lssh-periods}
% generated by make_tables.py from results/aeon-8585c7bf; do not edit
\begin{tabular}{lrcccccc}
\toprule
 & & \multicolumn{3}{c}{\texttt{AutoARIMA}} & \multicolumn{3}{c}{\texttt{AutoETS}} \\
\cmidrule(lr){3-5}\cmidrule(lr){6-8}
Sampling (period) & Series & whole & median & inv.\ OOS & whole & median & inv.\ OOS \\
\midrule
Daily (7) & 27 & 1.072 & 1.057 & 1.062 & 1.086 & 1.072 & 1.080 \\
Monthly (12) & 18 & 0.992 & 0.940 & 0.879 & 0.941 & 0.904 & 0.895 \\
Hourly (24) & 21 & 0.774 & 0.543 & 0.533 & 0.825 & 0.638 & 0.602 \\
Half-hourly (48) & 7 & 0.733 & 0.750 & 0.743 & 0.741 & 0.512 & 0.512 \\
Weekly (52) & 17 & 0.822 & 0.771 & 0.779 & 0.817 & 0.706 & 0.674 \\
10-minute (144) & 5 & 0.914 & 0.910 & 0.912 & 1.109 & 0.998 & 0.944 \\
Minutely (1440) & 5 & 1.492 & 0.718 & 0.746 & 1.192 & 0.641 & 0.728 \\
\bottomrule
\end{tabular}

\end{table}

\subsection{M4 long series}

Tables~\ref{tab:m4-main} and~\ref{tab:m4-schemes} give OWA on the 5{,}070
long M4 series. The gains are much smaller than on LSSH-99: the best schemes
reduce OWA from 0.920 to 0.892--0.898 for \texttt{AutoARIMA} and from 0.945 to
0.922--0.926 for \texttt{AutoETS}, about 2--3\%. They are concentrated where
LSSH-99 found them: hourly (\texttt{AutoARIMA} 0.692 to 0.578--0.648 for the
schemes in Table~\ref{tab:m4-main} other than the mean), monthly (0.917 to
0.850--0.876) and weekly series (0.878 to 0.829--0.860). Daily series,
58\% of the subset and close to random walks, show no effect for any method.
The M4 series are also shorter (median 1{,}430 against up to 10{,}000), and the
forecasts multi-step.

Very short windows are harmful on hourly series. Below 300 observations the M4
seasonality test, which considers lags up to $10\log_{10}n$, no longer reaches
lag 24, so the non-seasonal models are fitted without seasonal adjustment, and
\texttt{AutoETS} fits seasonal models to as few as 50 points. These windows ruin
the mean over windows (\texttt{AutoETS} OWA 1.027, worse than the whole series)
but not the median (0.926): the median does not depend on choosing the minimum
window well. Restricting the windows to at least 300 observations rescues the
mean on M4 (\texttt{AutoARIMA} 0.897) but hurts every scheme on LSSH-99
(\texttt{AutoETS} median 0.804 to 0.837), where short windows contribute to the
one-step forecasts.

The weighting schemes behave differently on the two benchmarks. Weights from
in-sample fit, which were among the best on LSSH-99, are no better than the
whole series on M4 for \texttt{AutoETS} (inverse in-sample MSE 0.943, relative
MAE 0.955, best in-sample window 0.960, Akaike-style 0.978). Weights from the
holdout error remain among the best (0.910 and 0.926 over all windows, 0.894 and
0.922 over windows of at least 300). The median of the better half of the
windows by out-of-sample error is among the four best schemes in all four
comparisons (LSSH-99 0.805 and 0.796; M4 0.898 and 0.924 for \texttt{AutoARIMA}
and \texttt{AutoETS}); the plain median and inverse out-of-sample MSE weights
are within 2\% of the best scheme in each comparison.

\begin{table}[t]
\centering\small
\caption{M4 series with at least 500 training observations: OWA relative to
Naive2 by frequency. Holdout selection uses windows of at least 300
observations; the combinations use all windows.}
\label{tab:m4-main}
\resizebox{\textwidth}{!}{% generated by make_tables.py from results/aeon-8585c7bf; do not edit
\begin{tabular}{lccccccc}
\toprule
 & Yearly & Quarterly & Monthly & Weekly & Daily & Hourly & All \\
\midrule
\multicolumn{8}{l}{\texttt{AutoARIMA}} \\
\quad Whole series & 1.019 & 0.998 & 0.917 & 0.878 & 1.007 & 0.692 & 0.920 \\
\quad Holdout selection ($w\geq300$) & 1.010 & 0.990 & 0.876 & 0.860 & 1.007 & 0.578 & 0.897 \\
\quad Mean & 1.005 & 0.971 & 0.858 & 0.829 & 1.001 & 1.221 & 0.975 \\
\quad Median & 1.011 & 0.973 & 0.855 & 0.843 & 1.002 & 0.648 & 0.905 \\
\quad Median of better half & 0.994 & 0.987 & 0.863 & 0.851 & 1.003 & 0.610 & 0.898 \\
\quad Inverse holdout MSE & 0.993 & 0.976 & 0.856 & 0.837 & 1.003 & 0.670 & 0.910 \\
\quad Inverse holdout MSE, $w\geq300$ & 1.002 & 0.993 & 0.874 & 0.844 & 1.002 & 0.600 & 0.894 \\
\quad Inverse in-sample MSE & 0.993 & 0.944 & 0.850 & 0.841 & 0.996 & 0.697 & 0.912 \\
\addlinespace
\multicolumn{8}{l}{\texttt{AutoETS}} \\
\quad Whole series & 0.990 & 0.939 & 0.869 & 0.963 & 1.002 & 0.788 & 0.945 \\
\quad Holdout selection ($w\geq300$) & 0.955 & 0.945 & 0.857 & 0.928 & 0.999 & 0.768 & 0.934 \\
\quad Mean & 0.988 & 0.913 & 0.844 & 0.934 & 0.999 & 1.697 & 1.027 \\
\quad Median & 0.992 & 0.934 & 0.836 & 0.932 & 1.003 & 0.750 & 0.926 \\
\quad Median of better half & 0.958 & 0.938 & 0.842 & 0.924 & 1.002 & 0.733 & 0.924 \\
\quad Inverse holdout MSE & 0.982 & 0.899 & 0.849 & 0.920 & 0.998 & 0.734 & 0.926 \\
\quad Inverse holdout MSE, $w\geq300$ & 0.982 & 0.937 & 0.850 & 0.931 & 0.999 & 0.719 & 0.922 \\
\quad Inverse in-sample MSE & 0.970 & 0.899 & 0.842 & 0.957 & 1.000 & 0.844 & 0.943 \\
\bottomrule
\end{tabular}
}
\end{table}

\begin{table}[t]
\centering\small
\caption{M4 long series: OWA of each combination scheme (all frequencies),
with the number of series with lower / higher MASE than the whole series.
$^*$: per-series MASE significantly different from the whole series in either
direction (Wilcoxon signed-rank, Bonferroni over 14 schemes). Bold: best per
forecaster.}
\label{tab:m4-schemes}
\resizebox{\textwidth}{!}{% generated by make_tables.py from results/aeon-8585c7bf; do not edit
\begin{tabular}{lcccc}
\toprule
Combination & \texttt{AutoARIMA} & \texttt{AutoETS} & \texttt{ARIMA}(1,0,1) & \texttt{ETS}(A,N,N) \\
\midrule
Whole series & 0.920 & 0.945 & 1.056 & 1.337 \\
Mean & 0.975$^*$ (2773/2273) & 1.027 (2684/2362) & 1.196$^*$ (2248/2822) & 1.336 (1860/1981) \\
Median & 0.905$^*$ (2627/1969) & 0.926$^*$ (2461/2097) & 1.064$^*$ (2647/2314) & 1.336 (1169/1121) \\
Trimmed mean (20\%) & 0.956$^*$ (2705/2240) & 0.927$^*$ (2636/2323) & 1.179 (2515/2555) & 1.336 (1684/1650) \\
Inverse in-sample MSE & 0.912$^*$ (2913/2139) & 0.943$^*$ (2793/2255) & 1.056 (2499/2571) & 1.336 (1983/2086) \\
Inverse in-sample relative MAE & 0.924$^*$ (2874/2172) & 0.955$^*$ (2715/2329) & 1.086 (2359/2711) & 1.336 (1859/1987) \\
Inverse squared in-sample relative MAE & 0.906$^*$ (2918/2128) & 0.939$^*$ (2795/2254) & 1.059 (2417/2653) & 1.336 (1860/1989) \\
Akaike-style weights ($n_0=100$) & 0.923$^*$ (2825/2199) & 0.978 (2597/2412) & 1.043$^*$ (2658/2407) & 1.338 (1931/2040) \\
Best in-sample window & 0.920$^*$ (2579/2191) & 0.960 (2248/2152) & 1.045 (2409/2369) & 1.338$^*$ (1281/1493) \\
Inverse out-of-sample MSE & 0.910$^*$ (2881/2170) & 0.926$^*$ (2796/2252) & 1.047 (2454/2616) & 1.336 (2035/2160) \\
Median of better half & 0.898$^*$ (2581/1978) & 0.924$^*$ (2373/2123) & 1.030$^*$ (2521/2304) & \textbf{1.336} (1116/1066) \\
Mean, $w\geq300$ & 0.897$^*$ (2929/2055) & 0.935$^*$ (2773/2210) & 1.035$^*$ (2671/2398) & 1.336 (1695/1618) \\
Median, $w\geq300$ & 0.899$^*$ (2605/1937) & 0.930$^*$ (2422/2071) & 1.035$^*$ (2685/2212) & 1.337 (1039/949) \\
Inverse in-sample MSE, $w\geq300$ & \textbf{0.892}$^*$ (2924/2069) & 0.938$^*$ (2743/2257) & \textbf{1.025}$^*$ (2749/2319) & 1.336 (1851/1799) \\
Inverse out-of-sample MSE, $w\geq300$ & 0.894$^*$ (2909/2087) & \textbf{0.922}$^*$ (2779/2210) & \textbf{1.025}$^*$ (2653/2415) & 1.336 (1776/1744) \\
\bottomrule
\end{tabular}
}
\end{table}

\subsection{Cost}
A window of 500--1{,}000 observations costs 2--4 times less to fit than the whole
series on LSSH-99 (median \texttt{AutoARIMA} fit 0.20--0.35\,s against 0.84\,s;
\texttt{AutoETS} 0.33--0.53\,s against 1.10\,s). Fitting every window costs
about 4 times the whole-series fit (median over series and origins 3.7--4.0,
interquartile range about 3.2--5.0), since most of the windows are short;
rolling selection
and out-of-sample weights additionally need the window forecasts at previous
origins, which in a rolling evaluation are available at no extra cost.

% =====================================================================
\section{A window-combining forecaster}
\label{sec:algorithm}

The results suggest a simple wrapper, Algorithm~\ref{alg:windower}. It turns any
forecaster into one that fits clones to a grid of recent windows and combines
their forecasts with the median, which needs no extra fits and is within 2\%
of the best scheme on both benchmarks; when out-of-sample errors per window are
available (from a holdout, or from previous origins in a rolling deployment),
the median of the better half of the windows is as robust and slightly more
accurate. Two details follow from
the experiments: the seasonal adjustment decision is taken once, on the whole
series, so that short windows do not switch it off; and the smallest window is
bounded below by a multiple of the seasonal period.

\begin{algorithm}[t]
\caption{\textsc{Windower}: fit and forecast with a combination of training windows}
\label{alg:windower}
\begin{algorithmic}[1]
\REQUIRE forecaster $F$, series $y_{1:t}$, horizon $h$, seasonal period $m$,
grid ratio $\rho$ (e.g.\ 1.5), minimum window $w_{\min}=\max(50, 3m)$,
combination $C\in\{\text{median}, \text{median of better half}\}$
\STATE decide seasonal adjustment on the whole series $y_{1:t}$ (if $F$ is non-seasonal)
\STATE $W \gets \{w_{\min}\rho^k : w_{\min}\rho^k < t\} \cup \{t\}$ \COMMENT{geometric grid plus the whole series}
\FOR{$w \in W$}
  \STATE fit a clone $F_w$ to $y_{t-w+1:t}$, with the seasonal adjustment of line 1
  \STATE $\hat y^{(w)}_{t+1:t+h} \gets$ forecast of $F_w$
\ENDFOR
\IF{$C$ is median}
  \RETURN $\mathrm{median}_{w\in W}\, \hat y^{(w)}_{t+1:t+h}$ (per step)
\ELSE
  \STATE $W' \gets$ the half of $W$ with the lowest $\mathrm{OOS}_w$ (holdout or past one-step errors)
  \RETURN $\mathrm{median}_{w\in W'}\, \hat y^{(w)}_{t+1:t+h}$
\ENDIF
\end{algorithmic}
\end{algorithm}

\todo{Implement as \texttt{Windower} in horizon, evaluate as an ordinary
forecaster on both benchmarks, and check it reproduces the offline results.
Cheaper variants: fewer windows; fixed model structure across windows;
incremental updates.}

% =====================================================================
\section{Discussion and limitations}

\begin{itemize}
\item \emph{When windows help.} The evidence is from real long series with
structure that changes over time. The variance cost of windowing on a stable
process is not measured here; simulated controls (stationary processes,
processes with breaks or drifting parameters) are needed to show when
windowing helps and when it costs.\todo{simulations}
\item \emph{Seasonal adjustment.} Windows below 300 change the seasonality
decision for hourly series under the M4 rule, so short-window results on
hourly data mix two effects.
\item \emph{Minutely wind power.} The largest gains (5 series of one dataset)
need inspection before being relied on.\todo{inspect}
\item \emph{Selection noise.} Choices made from 30 one-step errors, or from one
M4 holdout, are noisy; the oracle is optimistic, particularly on M4 with one
test period per series.
\item \emph{Multiple comparisons.} Many windows and schemes were compared on the
same test origins; following \citet{rossi12robust}, the headline results are
the schemes defined in advance (median, mean) rather than the best of many.
\item \emph{Other forecasters.} Only aeon's implementations were evaluated.
\end{itemize}

% =====================================================================
\section{Conclusion}
For long series, automatic ARIMA and ETS should not be fitted to all the data
by default. Fitting the forecaster to a geometric grid of recent windows and
taking the median of the forecasts reduced the error of \texttt{AutoARIMA} and
\texttt{AutoETS} by 11--13\% on LSSH-99 and their OWA by 2\% on the long M4
series, at about four times the cost of one fit. The gains are largest on
high-frequency seasonal series and absent on series close to a random walk.
Weighting windows by in-sample fit is not reliable; using out-of-sample error
to keep the better half of the windows helps slightly. Choosing a single window
is noisier than combining. \todo{revise after the wrapper evaluation and the
simulated controls.}

\bibliographystyle{plainnat}
\bibliography{refs}

\end{document}
